\documentclass[11pt,a4paper]{article}

\usepackage{kotex}
\usepackage[margin=24mm]{geometry}
\usepackage{array}
\usepackage{booktabs}
\usepackage{longtable}
\usepackage{enumitem}
\usepackage{hyperref}
\usepackage{xcolor}
\usepackage{listings}

\hypersetup{colorlinks=true,linkcolor=blue,urlcolor=blue}
\setlist{nosep}
\setlength{\parindent}{0pt}
\setlength{\parskip}{0.55em}
\renewcommand{\arraystretch}{1.25}

\lstset{
  basicstyle=\ttfamily\small,
  backgroundcolor=\color{gray!8},
  frame=single,
  breaklines=true,
  columns=fullflexible
}

\title{SW 암호모듈 기본설계서}
\author{교과 프로젝트}
\date{2026년 6월}

\begin{document}
\maketitle
\tableofcontents
\newpage

\section{문서 목적과 적용 범위}

본 문서는 기존 교과 코드인 ARIA, SHA3, HMAC-SHA3를 기반으로 구성한
소프트웨어 암호모듈의 기본설계를 설명한다. 모듈은 대칭키 암호, 해시,
메시지 인증, 난수 생성, 전통 공개키 키 합의, 양자내성 키 캡슐화 및
모듈 상태 통제를 하나의 공유 라이브러리 경계 안에 통합한다.

본 결과물은 수업의 최종 프로젝트를 위한 교육용 소프트웨어이다. KCMVP,
FIPS 140 또는 기타 공인 검증 절차에 제출되거나 공식 인증을 획득한
암호모듈이 아니다. 따라서 본 문서에서 ``지원'', ``시험 통과''라는 표현은
현재 저장소의 구현과 자체 시험 결과를 의미하며 공식 적합성 판정을
의미하지 않는다.

\section{설계 목표}

기본 설계 목표는 다음과 같다.

\begin{enumerate}
  \item 기존 교과 암호 코드를 재사용하고 공개 서비스 경계를 통일한다.
  \item 초기화 이전이나 오류 상태에서 암호 서비스가 실행되지 않도록 한다.
  \item 초기화 과정에서 무결성, 엔트로피 건전성 및 알고리즘 알려진 답 시험
        (KAT)을 수행한다.
  \item 난수에 의존하는 키 생성은 모듈 내부 Hash\_DRBG를 사용한다.
  \item P-256과 ML-KEM 키 생성 후 조건부 자가시험을 수행한다.
  \item 비밀 상태를 영점화하고 종료 후 서비스를 차단한다.
  \item 공유 라이브러리에서는 승인된 \texttt{KCMVP\_*} 심볼만 외부에 노출한다.
\end{enumerate}

\section{요구사항 충족 현황}

\begin{longtable}{p{0.22\textwidth}p{0.58\textwidth}p{0.12\textwidth}}
\toprule
요구사항 & 현재 구현 & 상태 \\
\midrule
\endhead
대칭키 암호 & ARIA-128/192/256, ECB/CBC/CTR 암호화 및 복호화 & 구현 \\
해시 & SHA3-224/256/384/512 & 구현 \\
MAC & HMAC-SHA3-224/256/384/512 & 구현 \\
RNG/DRBG & SHA3-256 기반 Hash\_DRBG, 생성/재시드/해제 & 구현 \\
엔트로피 & Linux \texttt{getrandom()} 전용 입력원 & 구현 \\
엔트로피 시험 & 시작 시 RCT와 APT 수행 & 구현 \\
ECC 키 합의 & P-256(secp256r1) ECDH, micro-ecc 사용 & 구현 \\
양자내성 KEM & PQClean clean 구현 기반 ML-KEM-768 & 구현 \\
시작 자가시험 & ARIA, SHA3, HMAC, DRBG, P-256, ML-KEM KAT & 구현 \\
무결성 시험 & 선택 소스 파일의 SHA3-256 생성 매니페스트 비교 & 구현 \\
조건부 시험 & P-256 및 ML-KEM 키 쌍 일관성 시험 & 구현 \\
상태 통제 & 비공개 FSM과 서비스 공통 인가 검사 & 구현 \\
영점화/종료 & DRBG 및 자가시험 상태 제거, 서비스 차단 & 구현 \\
공개 경계 & 공유 라이브러리와 버전 스크립트 기반 심볼 제한 & 구현 \\
\bottomrule
\end{longtable}

\section{전체 아키텍처}

시스템은 응용 프로그램, 공개 모듈 경계, 통제 기능, 암호 서비스 및 외부 참조
구현의 다섯 계층으로 구분한다.

\begin{lstlisting}
Application / Test Programs
  app/main.c, test/module_state_test.c, test/sha3_vs.c
                         |
                         v
Public API Boundary: include/kcmvp.h, lib/KCMVP_Crypto.c
  lifecycle | ARIA | SHA3/HMAC | RNG | ECDH | ML-KEM
                         |
          +--------------+---------------+
          | Control and assurance        |
          | FSM, self-test, integrity,    |
          | entropy health, zeroization   |
          +--------------+---------------+
                         |
          +--------------+---------------+
          | Cryptographic services       |
          | ARIA, SHA3, HMAC, Hash_DRBG,  |
          | P-256 adapter, ML-KEM adapter |
          +--------------+---------------+
                         |
          micro-ecc / PQClean reference implementations
\end{lstlisting}

응용 프로그램은 \texttt{include/kcmvp.h}만 포함하고
\texttt{build/libkcmvp\_crypto.so}에 연결한다. 원시 알고리즘 함수, 엔트로피
함수, DRBG 내부 함수와 외부 참조 프로젝트 심볼은 공개 서비스가 아니다.

\section{상태 기계 개요}

모듈 상태는 \texttt{module\_state.c}의 비공개 정적 변수로 유지된다. 외부
호출자는 상태를 조회할 수 있지만 직접 변경할 수 없다.

\begin{lstlisting}
LOAD
  -> PRE_SELFTEST
       -> NORMAL             모든 시작 시험 성공
       -> CRITICAL_ERROR     시작 시험 실패

NORMAL
  -> EXECUTION -> NORMAL     일반 암호 서비스
  -> COND_SELFTEST -> NORMAL 조건부 자가시험 성공
  -> CRITICAL_ERROR          치명 오류 또는 zeroize
  -> EXIT                    정상 shutdown

EXECUTION / COND_SELFTEST
  -> CRITICAL_ERROR          서비스 또는 조건부 시험의 치명 실패

CRITICAL_ERROR
  -> EXIT                    shutdown만 허용
\end{lstlisting}

\texttt{DEGRADED}와 \texttt{NORMAL\_ERROR}는 공개 열거형에 예약되어 있으나
현재 구현 경로에서는 사용하지 않는다. 초기화 전 서비스는
\texttt{KCMVP\_ERROR\_NOT\_INITIALIZED}, 정상 상태가 아닌 경우에는 주로
\texttt{KCMVP\_ERROR\_INVALID\_STATE}를 반환한다.

\section{지원 알고리즘과 서비스}

\subsection{ARIA}

ARIA는 128, 192, 256비트 키를 지원한다. 공개 모드는 ECB, CBC, CTR이다.
ECB와 CBC 입력 길이는 16바이트의 양의 배수여야 하며 패딩은 호출자의
책임이다. CTR은 마지막 부분 블록을 포함한 임의의 양의 바이트 길이를
처리한다. CBC 및 CTR을 포함하여 암호화와 복호화를 모두 제공한다.

\subsection{SHA3와 HMAC-SHA3}

SHA3-224/256/384/512 단일 호출 해시와 동일 출력 크기의 HMAC-SHA3를
제공한다. SHA3 내부 스트리밍 상태는 명시적 문맥 구조체에 저장되므로 서로
독립적인 해시 계산이 전역 Keccak 상태를 공유하지 않는다.

\subsection{Hash\_DRBG}

SHA3-256을 사용하는 프로젝트용 Hash\_DRBG이다. 내부 상태는 440비트의
\texttt{V}, \texttt{C}, 재시드 카운터 및 인스턴스 상태로 구성된다. 한 번의
생성 요청은 최대 65,536바이트이며 재시드 간격은 $2^{48}$ 요청이다.
초기화와 재시드 엔트로피는 \texttt{entropy\_get()}을 통해 얻는다.

\subsection{P-256 ECDH}

BSD 2-Clause 라이선스의 micro-ecc 참조 코드를 사용한다. 개인키는 모듈
Hash\_DRBG로 생성한다. 공개키 형식은 SEC1 접두 바이트가 없는 64바이트
$X \parallel Y$이고, 공유 비밀은 32바이트 x 좌표이다. 상대 공개키는 곡선
유효성 검사를 거친다.

\subsection{ML-KEM-768}

PQClean의 portable clean ML-KEM-768 구현을 사용한다. 공개키는 1,184바이트,
비밀키는 2,400바이트, 암호문은 1,088바이트, 공유 비밀은 32바이트이다.
키 생성과 캡슐화 난수는 모듈 Hash\_DRBG로 연결된다. PQClean 소스는
CC0/public domain 조건이며 저장소에 고정된 커밋 정보가 기록되어 있다.

\section{초기화 및 시험 전략}

\texttt{KCMVP\_Initialize()}의 논리적 순서는 다음과 같다.

\begin{enumerate}
  \item \texttt{LOAD}에서 \texttt{PRE\_SELFTEST}로 전이한다.
  \item 생성된 SHA3-256 매니페스트로 선택 모듈 파일의 무결성을 확인한다.
  \item Linux \texttt{getrandom()}으로 512바이트를 수집한다.
  \item Repetition Count Test(RCT)와 Adaptive Proportion Test(APT)를 수행한다.
  \item ARIA-128 암호화/복호화 KAT를 수행한다.
  \item SHA3-256 및 HMAC-SHA3 KAT를 수행한다.
  \item 결정론적 Hash\_DRBG KAT를 수행한다.
  \item 결정론적 P-256 공개키/ECDH KAT를 수행한다.
  \item 결정론적 ML-KEM-768 키 생성/캡슐화 KAT를 수행한다.
  \item 새 엔트로피로 운영용 Hash\_DRBG를 인스턴스화한다.
  \item 모든 단계 성공 시 \texttt{NORMAL}로 전이한다.
\end{enumerate}

각 알고리즘 통과 플래그는 해당 KAT 성공 후에만 설정한다. 어느 단계든
실패하면 시험 및 DRBG 상태를 제거하고 \texttt{CRITICAL\_ERROR}로 전이하여
암호 서비스를 차단한다.

운영 중 P-256과 ML-KEM 키 쌍 생성은 결과를 반환하기 전에 조건부 자가시험을
수행한다. P-256은 상호 ECDH 비밀을 비교하고, ML-KEM은 결정론적 캡슐화 후
복호화한 공유 비밀을 비교한다. 실패 시 생성 키를 영점화하고 모듈을
\texttt{CRITICAL\_ERROR}로 전이한다.

\section{실행 및 검증 방법}

Linux, GCC, GNU Make, binutils가 필요하다. 무결성 매니페스트와 벡터 경로가
프로젝트 루트 기준이므로 명령은 저장소 루트에서 실행한다.

\begin{lstlisting}[language=bash]
make clean
make

# 번호 메뉴 또는 전체 공개 서비스 시연
./build/main
./build/main --all

# 전체 시험과 공개 심볼 감사
make test
make audit-exports

# clean build + tests + export audit
make clean && make audit
\end{lstlisting}

주요 시험은 상태 및 장애 주입 시험, 3,342개 ARIA 벡터, 1,266개
SHA3/HMAC-SHA3 벡터, P-256 상호 합의, ML-KEM 캡슐화/복호화, 무결성 실패,
엔트로피 실패, KAT 실패 및 조건부 시험 실패를 포함한다.

\section{제약 및 한계}

\begin{itemize}
  \item 본 모듈은 공식 KCMVP/FIPS 검증 모듈이 아니다.
  \item ECB/CBC는 패딩을 제공하지 않는다. CFB64/OFB는 공개 서비스가 아니다.
  \item 무결성 시험은 선택한 소스 파일의 해시를 비교하며 서명된 코드 로딩이나
        메모리에 적재된 실행 페이지의 검증을 제공하지 않는다.
  \item 무결성 경로는 프로젝트 루트에서 실행하는 현재 빌드 구조에 의존한다.
  \item 엔트로피 시험 임계값은 프로젝트 설정이며 실제 엔트로피 원천의
        최소 엔트로피 평가에 의해 정당화된 인증 파라미터가 아니다.
  \item Hash\_DRBG 구현은 교육용 프로젝트 코드이며 독립 표준 검증을 받지 않았다.
  \item ECDH 출력은 원시 공유 비밀이다. 응용 키로 사용하려면 적절한 KDF가 필요하다.
  \item 호출자에게 반환된 개인키와 공유 비밀의 수명 및 영점화는 호출자 책임이다.
  \item PQClean은 고정 커밋을 사용하며 해당 상위 프로젝트의 유지보수 상태 변화에
        주의해야 한다.
\end{itemize}

\end{document}
